% !TEX root = ../diplomarbeit.tex

\section{\emph{Beispiel:} Hardware\authorA}

% Hier beschreibt Autor A seinen Teil des Prototyps.
% Beschreibung der Umsetzung und durchgeführten Tests.

% --------------------------------------------------------
% Entwicklung und Funktionalitäten des Prototyps
% --------------------------------------------------------
\subsection{Beschreibung des Prototyps}

\emph{Hier wird der entwickelte Prototyp mit seinen Funktionen, der Benutzeroberfläche und den umgesetzten Kernfunktionalitäten vorgestellt. Screenshots, Ablaufdiagramme und Beschreibungen der Benutzerinteraktion dokumentieren den aktuellen Entwicklungsstand.}

\placeholderbox{Beispiel}{%
Der Prototyp wurde als Web-Anwendung mit React (Frontend) und Node.js (Backend) implementiert. Die Benutzeroberfläche besteht aus drei Hauptkomponenten: (1) ein Eingabebereich für Konfigurationsparameter, (2) eine Visualisierungskomponente für Echtzeitdaten und (3) ein Export-Panel für Ergebnisse. Die Anwendung arbeitet vollständig offline und speichert Daten lokal in IndexedDB. Die Kernfunktionalitäten wurden durch 42 Unit-Tests validiert.%
    %
    \newline\newline
    ...
}

% --------------------------------------------------------
% Gewählte Testverfähren und Ergebnisse
% --------------------------------------------------------
\subsection{Tests und Evaluierung}

\emph{Dieser Abschnitt beschreibt die durchgeführten Tests (Unit-Tests, Integrationstests, Benutzertests) sowie deren Ergebnisse. Die Evaluierung bewertet, ob die definierten Anforderungen erfüllt wurden und identifiziert Verbesserungspotenziale.}

\placeholderbox{Beispiel}{%
Folgende Testabdeckung wurde durchgeführt: (1) Unit-Tests für alle kritischen Funktionen mit 98,1\,\% Erfolgsquote, (2) Integrationstests mit realen Eingabedaten zeigten eine durchschnittliche Verarbeitungszeit von 2,3 Sekunden, (3) Usability-Tests mit 8 Testnutzern erreichten einen SUS-Wert von 72 Punkten. Browser-Kompatibilität wurde auf Chrome, Firefox und Safari getestet. Ein kritischer Bug in der Datenvalidierung wurde identifiziert und behoben.%
    %
    \newline\newline
    ...
}

% --------------------------------------------------------
% Erreichte Ergebnisse und Metriken
% --------------------------------------------------------
\subsection{Ergebnisse}

\emph{Dieser Abschnitt fasst die wichtigsten Ergebnisse der Evaluierung zusammen. Die Resultate werden anhand quantitativer Metriken sowie qualitativer Beobachtungen dargestellt und in Bezug zu den in der Zielsetzung definierten Anforderungen bewertet.}

\placeholderbox{Beispiel}{
    Bei der Durchführung von insgesamt 42 Unit-Tests erreichte der Prototyp eine Erfolgsquote von 98,1\,\%. Die durchschnittliche Antwortzeit bei typischen Anwendungsfällen betrug 2,3 Sekunden, was die Anforderung von unter 5 Sekunden deutlich erfüllt. %
    %
    \newline\newline
    ...
}


\section{\emph{Beispiel:} Frontend\authorA}

\subsection{Beschreibung des Prototyps}

\ldots

\subsection{Tests und Evaluierung}

\ldots

\subsection{Ergebnisse}

\ldots